\section{Article Retrieval and Knowledge Extraction}

The recommendation process in MLguide depends on a database of 
research articles, where problem descriptions are matched 
through semantic retrieval. This article base is built and 
maintained by a separate article retrieval pipeline, which is the 
subject of this section.

Article abstracts are used as the textual basis for each article. 
Text mining of scientific literature has mostly been carried out on 
abstracts, both because they are widely available and because they 
consist of short sentences presenting only the main findings of a 
study \parencite{westergaard2018comprehensive_quantitative_comparison_text_mining_15_million_full_text_articles_versus_corresponding_abstracts}. 
Full-text articles contain richer information, but also include 
tables, references, and discussion sections that make them less 
consistent as a basis for comparison across a large corpus 
\parencite{westergaard2018comprehensive_quantitative_comparison_text_mining_15_million_full_text_articles_versus_corresponding_abstracts}. 
Abstracts therefore provide a consistently available representation 
of each article's core content.

Scopus was selected as the source of research articles. Scopus and 
Web of Science are the two main databases for citation data 
\parencite{mongeon2016journal_coverage_web_science_scopus}, and of the two, 
Scopus has the broader journal coverage 
\parencite{mongeon2016journal_coverage_web_science_scopus}. Scopus also 
provides an API \parencite{scopus_api}, which allows article metadata and abstracts to be 
retrieved automatically. This makes it possible to collect articles 
in a repeatable way and to expand the article database over time.


\subsection{Data Collection}
\label{sec:article-retrieval-data-collection}
 
The article retrieval pipeline builds on the data collection and
classification work from the pre-master project
\parencite{bergsto2026mlselection}. The same query-based search
strategy is reused here, and the pipeline is adapted so that new
articles can be added to the MLguide databases over time.
 
Articles are collected from Scopus through a set of queries, each
targeting machine learning applications in a specific domain or
application area. Each query is defined in a YAML configuration
file with an identifier and a Scopus search string. The full set
of queries is listed in Appendix~\ref{app:queries}, and the search
strategy behind them is described in detail in the pre-master project
\parencite{bergsto2026mlselection}. When the pipeline runs, each
query is sent to the Scopus API \parencite{scopus_api} through the
Python package pybliometrics \parencite{pybliometrics_rose_kitchin_2019}, which
provides a stable interface for retrieving data from Scopus. For
each article, the DOI, Scopus EID, title, and abstract are stored,
together with the identifier of the query that retrieved it. The
DOI and EID are used to identify articles and to skip articles
that are already present in the database, to avoid duplicates.

\subsection{Adaptations for MLguide}
\label{sec:adaptations-for-mlguide}

The pre-master project classified each article along four
dimensions: ML paradigm, product lifecycle (PLC) phase,
application area, and the ML methods mentioned in the abstract.
The pipeline in MLguide keeps only ML paradigm and ML methods,
as these are the dimensions used in the recommendation process.
ML paradigm is used as the ML area filter, and ML methods link
each article to methods in the knowledge graph.

PLC phase and application area were removed because they were
not reliable enough to act as filters. As discussed in the
pre-master project \parencite{bergsto2026mlselection}, single-phase
PLC labels are only an approximation, since many studies span
several lifecycle phases, and the application areas, represented by clusters, 
did show limited natural separation. Using either as a filter would risk excluding
relevant articles on uncertain grounds. Instead, the system relies
on semantic similarity between the problem description and the
article abstracts, as described in
Section~\ref{sec:recommendation-process}.

\subsection{Knowledge Extraction}
\label{sec:article-retrieval-knowledge-extraction}

After articles are retrieved, two types of information are
extracted from each abstract: the ML methods mentioned in the
text, and the ML paradigm the article belongs to. Both steps work
on a cleaned version of the abstract, where copyright statements,
URLs, and similar noise are removed so that matching is more
reliable.

ML methods are extracted using a dictionary of known methods. Each
method is linked to a set of common phrasings, and the cleaned
abstract is searched for these phrasings using word-boundary
matching. A method is recorded for an article if any of its
phrasings is found. When a specific neural network variant is
already detected, the generic term \texttt{neural network} is
removed, so that articles are linked to the most specific method
mentioned.

ML paradigm is assigned with a keyword-based approach. Separate
keyword lists are defined for supervised, unsupervised, and
reinforcement learning, and the abstract is matched against each
list. An article can be assigned more than one paradigm if terms
from several lists are found. Articles with no matches are labelled
as unknown. This keyword approach is used because it is simple to
run for every batch and does not depend on a trained model.

\subsection{Database Upload}
\label{sec:article-retrieval-database-upload}

The final step stores the processed articles in the two MLguide
databases. Abstracts and embeddings are uploaded to PostgreSQL,
and the links between articles and ML methods are uploaded to
GraphDB.

For PostgreSQL, each new article is stored with its cleaned
abstract, and embeddings are computed for both the title and the
abstract. The embeddings are produced with the
\texttt{all-mpnet-base-v2} sentence-transformer model, which maps
text to 768-dimensional vectors. The same model is used in the
recommendation process, so that problem descriptions and stored
articles are placed in the same embedding space. Articles already
present in the database are skipped, so only new articles are
added.

For GraphDB, the extracted methods and paradigms are converted to
RDF triples in Turtle format and uploaded to the knowledge graph.
Each article is linked to the ML methods it mentions and tagged
with its ML area. These links allow the recommendation process to
move from articles to methods when retrieving candidates, as
described in Section~\ref{sec:recommendation-process}.

\label{sec:article-retrieval-and-knowledge-extraction}
\begin{figure}[H]
    \centering
    \includegraphics[width=0.9\textwidth]{Figures/06-implementation/article-retrieval-flow.pdf}
    \caption[Article retrieval flow]{Flowchart of the article retrieval and knowledge extraction process in MLguide.}
    \label{fig:article-retrieval-flow}
\end{figure}


\textbf{Bruke figuren mer aktivt}
\textbf{Nevne på slutten antall artikler som kom initielt og at det ble lagt til artikler i en til runde på en bestemt dato.}